\begin{Definition}
\label{def:symmetric-operator}
An operator $(\mathcal{D}(A),A)$ from $H$ to $H$ with dense domain $\mathcal{D}(A)$ is said to be \textbf{symmetric} if $(\mathcal{D}(A^*),A^*)$ is an extension of $(\mathcal{D}(A),A)$, or equivalently if
\begin{equation}\label{eq:symmetric-operator-pairing-identity}
(Au,v)=(u,Av)\quad\text{for every }u,v\in\mathcal{D}(A).
\end{equation}

\end{Definition}

\begin{Definition}
\label{def:self-adjoint-operator}
An operator is said to be \textbf{self-adjoint} if
\begin{equation}\label{eq:self-adjoint-operator-domain-action-equality}
(\mathcal{D}(A),A)=(\mathcal{D}(A^*),A^*).
\end{equation}

\end{Definition}

The following lemma is a direct consequence of the definitions.

\begin{Lemma}
\label{lem:self-adjoint-symmetric-domain-equality}
An operator $(\mathcal{D}(A),A)$ is self-adjoint if and only if it is symmetric and $\mathcal{D}(A)=\mathcal{D}(A^*)$.
\end{Lemma}

For a bounded operator, symmetry and self-adjointness are the same thing.

\begin{Theorem}
\label{thm:bounded-symmetric-operator-self-adjoint}
An operator $A\in\mathcal{L}(H)$ is self-adjoint if and only if it is symmetric.
\end{Theorem}

\begingroup
\postdisplaypenalty=10000
\renewcommand{\endProof}{\par\nobreak\hfill$\blacksquare$\par\vspace{0.1cm}}
\begin{Proof}
By definition self-adjointness implies symmetry. Furthermore, if $\mathcal{D}(A)=H$ and $A$ is symmetric, then by \eqref{eq:symmetric-operator-pairing-identity}, $\mathcal{D}(A^*)=\mathcal{D}(A)=H$.
\end{Proof}
\endgroup

As we shall see in Problem~\ref{ex:second-derivative-symmetric-not-self-adjoint}, there are examples of unbounded symmetric operators that are not self-adjoint.

\Needspace{4\baselineskip}
\begin{Example}
\label{ex:differentiation-hilbert-adjoint-complete-calculation}
We now give an example of the computation of the adjoint of an unbounded operator. Let $H=L^2(0,1)$ and let\[\mathcal{D}(D)=\{u\in H^1(0,1)\mid u(0)=0\}.\]Here the boundary condition is taken in the sense of trace. Define the differentiation operator $(\mathcal{D}(D),D)$ from $H$ to $H$ by\[D(u):=u'\]for all $u\in\mathcal{D}(D)$. We begin computing the adjoint by doing a formal calculation which gives us a ``guess'' as to the identity of the adjoint. We then do a rigorous proof that the guess was right.

Both the formal calculation and the rigorous proof are based on the identity
\begin{equation}\label{eq:differentiation-adjoint-defining-pairing}
(v,Du)=(D^*v,u)\quad\text{for all }u\in\mathcal{D}(D),\ v\in\mathcal{D}(D^*).
\end{equation}
For the formal calculation we begin on the left of \eqref{eq:differentiation-adjoint-defining-pairing} and proceed to integrate by parts (even though we do not yet know anything about $v$ other than that it is in $L^2(0,1)$). We get\[\begin{aligned}
(v,Du)&=\int_0^1u'(x)v(x)\,\mathrm{d}x\\
&=-\int_0^1u(x)v'(x)\,\mathrm{d}x+u(1)v(1)-u(0)v(0)\\
&=(u,-v')+u(1)v(1).
\end{aligned}\]In the last equality we used the fact that $u(0)=0$. Examining \eqref{eq:differentiation-adjoint-defining-pairing} (and requiring that our formal calculations make sense) leads us to guess that the adjoint is given by $(\mathcal{D}(B),B)$ where
\begin{equation}\label{eq:differentiation-adjoint-candidate-domain-action}
\mathcal{D}(B):=\{v\in H^1(0,1)\mid v(1)=0\},\quad B(v)=-v'.
\end{equation}
This guess is correct, but all we have shown at this time is that the adjoint is an \emph{extension} of the operator $(\mathcal{D}(B),B)$. To prove that $(\mathcal{D}(D^*),D^*)=(\mathcal{D}(B),B)$, we let $D^*v:=f$ (making no assumptions other than that it is in $L^2(0,1)$) and define
\begin{equation}\label{eq:differentiation-adjoint-antiderivative}
F(x)=\int_x^1f(s)\,\mathrm{d}s.
\end{equation}
Of course, we note from \eqref{eq:differentiation-adjoint-candidate-domain-action} that $F\in\mathcal{D}(B)$ and that $B(F)=f$. We now work from the \emph{right} of \eqref{eq:differentiation-adjoint-defining-pairing} and get that for every $v\in\mathcal{D}(D^*)$ and $u\in\mathcal{D}(D)$ we have\[\begin{aligned}
(D^*v,u)&=(f,u)\\
&=\int_0^1u(x)f(x)\,\mathrm{d}x\\
&=-\int_0^1u(x)F'(x)\,\mathrm{d}x\\
&=\int_0^1u'(x)F(x)\,\mathrm{d}x\\
&=(F,Du)\\
&=(v,Du).
\end{aligned}\]Since this is true for all $u\in\mathcal{D}(D)$ we must have $F-v\in\mathcal{R}(D)^\perp$. However, $\mathcal{R}(D)=L^2(0,1)$, so $F=v$. Hence, the assertion is proved.
\end{Example}

\Needspace{4\baselineskip}
\begin{Example}
\label{ex:general-differential-hilbert-formal-adjoint}
For a more general partial differential operator, such as those defined in Example~\ref{ex:differential-operator-boundary-domains}, calculation of the adjoint is complicated by the fact that we cannot yet calculate the closure of the operator. However, without closing the operator, one can show that the Hilbert adjoint is an \emph{extension} of the formal adjoint defined in Definition~\ref{lem:distribution-convolution-continuity}.
\end{Example}

As we shall see below, self-adjoint operators have many properties that are useful in application. However, as the example above indicates, it is often difficult to determine whether an unbounded operator is self-adjoint. (We must first close the operator, and as we have seen, finding the domain of the closure is often nontrivial.) Fortunately, we can define a related property called \emph{essential self-adjointness} for which there is a relatively easy test.
